\section{A fully constructed infinite nonlinear system}
\label{sec:lattice}

We now eliminate every abstract domain, evaluator and tail premise in
one example. The equation is
\[
 x_i=a_i+\frac{x_i^2}{8}+\frac{x_{i+1}}8,
 \qquad a_i=\frac1{4(i+1)},\qquad 0\le x_i\le1
 \quad(i\in\N).
 \label{eq:lattice-system}\tag{L1}
\]
This is a genuinely infinite dependency: computing a coordinate by
recursion on its successors would not terminate. The construction
below instead proves a bounded amount of successor information
sufficient for any specified precision.

\subsection{Finite models and their certificates}

For \(n\ge1\), let \(D_n=([0,1]\cap\Q)^n\), with
\[
 d_n(x,y)=\max_{i<n}2^{-i}|x_i-y_i|.
\]
These are weighted maximum distances with weights \(w_i=2^i\).
Let \(P_n\) delete the last coordinate. Prefix restrictions are
nonexpansive and compose exactly.

Define \(F_n\) by
\[
 (F_nx)_i=a_i+\frac{x_i^2}{8}+\frac{x_{i+1}}8,
 \quad i<n,
 \qquad x_n:=0.
 \label{eq:lattice-finite}\tag{L2}
\]
All coordinates are nonnegative, and
\((F_nx)_i\le1/4+1/8+1/8=1/2\). Thus the map preserves the box,
with room for downward rounding.

For \(u,v\in[0,1]\),
\(\abs{u^2-v^2}=\abs{u-v}(u+v)\le2\abs{u-v}\).
Consequently
\[
 |(F_nx)_i-(F_ny)_i|
 \le\frac14|x_i-y_i|+\frac18|x_{i+1}-y_{i+1}|.
\]
The absent last-coordinate difference is zero. Multiplication by
\(2^{-i}\) bounds the two terms by \(d_n(x,y)/4\) and
\(d_n(x,y)/4\), respectively. Hence
\[
 d_n(F_nx,F_ny)\le\tfrac12d_n(x,y).             \tag{L3}
\]
The gain matrix has diagonal entries \(1/4\), superdiagonal entries
\(1/8\), and zeros elsewhere. With \(w_i=2^i\), every nonterminal
weighted row equals \(w_i/2\). This is a literal finite instance of
the matrix certificate.

The vectors \(F_nP_nz\) and \(P_nF_{n+1}z\) agree at every
coordinate except \(n-1\); there the difference is \(z_n/8\).
Thus
\[
 d_n(F_nP_nz,P_nF_{n+1}z)
 \le2^{-(n-1)}/8=2^{-n-2}.                     \tag{L4}
\]
This is the required commutation defect. It comes from the boundary
truncation, not from a numerical roundoff estimate.

\subsection{A bounded dyadic implementation}

For a nonnegative rational \(u\), define
\[
 \lfloor u\rfloor_p=2^{-p}\lfloor2^pu\rfloor.
\]
Integer division computes it, and
\(0\le u-\lfloor u\rfloor_p<2^{-p}\). Use
\emph{synchronous} iterations: all new coordinates are computed from
the previous vector. Set
\begin{align*}
 p_n&=n+8, & t_n&=n+4,\\
 v_{n,0}&=0, &
 v_{n,t+1,i}&=\left\lfloor(F_nv_{n,t})_i\right\rfloor_{p_n},\\
 y_n&=v_{n,t_n}.&&
\end{align*}
The rounded states lie in \([0,1/2]^n\). Each step has weighted
error at most \(h_n=2^{-p_n}\), since every factor \(2^{-i}\le1\).

Let \(z_{n,t}=F_n^t(0)\) be the exact rational iterates. These are
used for the proof; the reference implementation need not construct
their large denominators. With
\(\delta_t=d_n(v_{n,t},z_{n,t})\), (L3) gives
\[
 \delta_0=0,\qquad
 \delta_{t+1}\le\tfrac12\delta_t+h_n,
 \qquad \delta_t\le2h_n.                       \tag{L5}
\]
The last bound is a single quantifier-free induction.

The initial exact increment has distance
\(d_n(0,F_n0)=1/4\). Repeated use of (L3) bounds the exact residual
after \(t\) steps by \(2^{-t-2}\). Therefore
\begin{align*}
 d_n(y_n,F_ny_n)
 &\le d_n(y_n,z_{n,t_n})
       +d_n(z_{n,t_n},F_nz_{n,t_n})
       +d_n(F_nz_{n,t_n},F_ny_n)\\
 &\le3h_n+2^{-t_n-2}
   =7\,2^{-n-8}
   \le2^{-n-4}.
\end{align*}
This is a proved residual budget, not a stopping criterion based on
observed agreement of successive decimals.

\subsection{Coherence and the actual coordinate formula}

Apply \cref{thm:commutation} with \(q_n=1/2\), \(L_n=1\),
\(\eta_n=2^{-n-2}\), and \(r_n=2^{-n-4}\). It yields
\[
 d_n(y_n,P_ny_{n+1})
 \le2(2^{-n-4}+2^{-n-2}+2^{-n-5})
 =\frac{11}{16}2^{-n}<2^{-n}.                  \tag{L6}
\]
The finite geometric sum and \cref{thm:finite-refinement} imply
\[
 d_n(y_n,P_{n,m}y_m)\le2^{1-n}\quad(m\ge n).
 \tag{L7}
\]
Coordinate \(i\) is an observer at every level \(n>i\), with
Lipschitz bound \(2^i\). Define the output by the explicit formula
\[
 \boxed{X(i,k)=(y_{i+k+3})_i.}                 \label{eq:lattice-name}\tag{L8}
\]

\begin{theorem}[Uniform solution sampler]
\label{thm:lattice-name}
The function (L8) is primitive recursive and takes values in \([0,1]\).
For every \(i,k,l\in\N\),
\begin{align}
 |X(i,k)-X(i,l)|&\le2^{-\min(k,l)-2},          \tag{L9}\\
 \left|X(i,k)-a_i-\frac{X(i,k)^2}{8}
                      -\frac{X(i+1,k)}8\right|
 &\le2^{-k-4}.                               \tag{L10}
\end{align}
Thus it represents a solution of (L1) by entirely arithmetic contracts.
\end{theorem}
\begin{proof}
Every operation in (L8) is a bounded list operation, integer division,
or a PR iteration with the explicit bounds \(n=i+k+3\),
\(t=n+4\), \(p=n+8\). Membership in the box has already been
proved. If \(k\le l\), (L7) gives
\[
 |X(i,k)-X(i,l)|\le2^i2^{1-(i+k+3)}=2^{-k-2}.
\]
The reverse order gives (L9).

For (L10), put \(n=i+k+3\). In particular \(i+1<n\), so
the \(i\)-th equation of \(F_n\) uses the actual next coordinate,
not the zero boundary. Its residual contributes at most
\(2^ir_n\). The value \(X(i+1,k)\), however, is taken at level
\(n+1\). Replacing \((y_n)_{i+1}\) by \((y_{n+1})_{i+1}\)
contributes at most \(2^{i+1}2^{-n}/8\), by (L6). Therefore
\[
 \left|X(i,k)-a_i-\frac{X(i,k)^2}{8}-\frac{X(i+1,k)}8\right|
 \le2^{i-n-4}+2^{i-n-2}
 =5\,2^{-k-7}<2^{-k-4}.
\]
This also explains why the difference between successive model
dimensions must be included in the equation proof.
\end{proof}

\subsection{Uniqueness in the specified box}

Call a candidate admissible if it is a regular rational array \(U(i,p)\)
with \(0\le U(i,p)\le1\), and its same-precision residual satisfies
\[
 \left|U(i,p)-a_i-\frac{U(i,p)^2}{8}
                     -\frac{U(i+1,p)}8\right|\le2^{-p}
 \quad(i,p\in\N).                              \tag{L11}
\]
This is one useful normalized representation. Ordinary regular
solution names with a known residual modulus can be reindexed to
this format. No algorithm for extracting such a modulus from an
unspecified external function is assumed.

\begin{theorem}[Finite uniqueness certificate]
\label{thm:lattice-unique}
If \(U,V\) are admissible, then, for every \(i,k\), putting
\(p=i+k+5\) gives
\[
 |U(i,p)-V(i,p)|\le2^{-k-2}.                    \tag{L12}
\]
Consequently their represented coordinates are equal. In particular,
(L8) is the unique represented solution in this class.
\end{theorem}
\begin{proof}
Let \(N=i+k+3\), and form the first \(N\) coordinates of
\(U(-,p)\) and \(V(-,p)\), denoted \(u,v\). For each of the
first \(N-1\) equations their finite-map residual is at most
\(2^{-p}\). The last equation has the additional omitted boundary
coordinate, contributing at most \(2^{-N-2}\) in weighted distance.
Thus both finite-map residuals are at most
\(2^{-p}+2^{-N-2}\). The triangle inequality and (L3) give
\[
 d_N(u,v)\le2(2^{-p}+2^{-N-2})+\tfrac12d_N(u,v),
\]
and hence \(d_N(u,v)\le4\,2^{-p}+2^{-N}\). Multiplication by
\(2^i\) yields
\[
 |U(i,p)-V(i,p)|
 \le2^{i+2-p}+2^{i-N}
 =2^{-k-3}+2^{-k-3}=2^{-k-2}.
\]
Only coordinates through \(N\) and one finite precision are queried.
For an entirely explicit equality bound at a fixed shift, apply
(L12) with target index \(k+2\), so the common work precision is
\(p=i+k+7\). Regularity of each candidate, between \(k+2\)
and \(p\), then gives
\[
 |U(i,k+2)-V(i,k+2)|
 \le 2\bigl(2^{-k-2}+2^{-i-k-7}\bigr)+2^{-k-4}
 \le\frac{37}{64}2^{-k}<2^{-k}.
\]
Thus equality itself has a supplied finite precision schedule. The
constructed sampler satisfies the stronger residual (L10), so it is
admissible.
\end{proof}

\subsection{The box is mathematically necessary}

Uniqueness among \emph{all bounded} solutions would be false. The
following independent finite construction produces another solution,
with every coordinate in \([6,8]\).

\paragraph{A completely finite upper-branch construction.}
For $z\in[6,8]\cap\Q$, put
$R_i(z)=16-z-2/(i+1)\in[6,10]$.
Starting with $[2,4]$, perform $p+4$ rational bisections, retaining an
interval $[l,u]$ with $l^2\le R_i(z)\le u^2$; at each step test the
square of the midpoint.  Define $b_i(z,p)$ to be the final midpoint and
$E_i(z,p)=4+b_i(z,p)$.  The final interval width is $2^{-p-3}$.
Since $2\le l\le b_i\le u\le4$,
\[
 |b_i(z,p)^2-R_i(z)|\le 8\,(u-l)/2=2^{-p-1}.
\]
The exact identities for differences of squares give, entirely over
rationals,
\[
 \begin{split}
 6\le E_i(z,p)&\le8,\\
 |E_i(z,p)-E_i(z',q)|
 &\le\tfrac14|z-z'|+\tfrac18(2^{-p}+2^{-q}),\\
 \left|E_i(z,p)-a_i-E_i(z,p)^2/8-z/8\right|
 &\le2^{-p}/16.
 \end{split}
\]
For the second inequality, factor the difference of the two midpoint
squares and use that their sum is at least four.  For the third, multiply
the expression by eight; it becomes $R_i(z)-b_i(z,p)^2$.

Define the finite backward chain by
$x_N^{N,p}=7$ and $x_j^{N,p}=E_j(x_{j+1}^{N,p},p)$ for $j<N$.
For $M\ge N$ and $i<N$, its cross-precision estimate is
\[
 |x_i^{N,p}-x_i^{M,q}|
 \le4^{-(N-i)}+\tfrac16(2^{-p}+2^{-q}).
\]
Indeed the terminal discrepancy at coordinate $N$ is at most one, and
finite iteration of the preceding contraction inequality gives the first
term and a geometric sum bounded by
$(2^{-p}+2^{-q})(1/8)(4/3)$.

Now define the actual rational sampler
\[
 Y(i,k)=x_i^{\,i+k+3,\,k+4}.
\]
It uses exactly $k+3$ backward applications, each with $k+8$ bisections.
For $k\le l$, the cross-precision estimate gives
\[
 |Y(i,k)-Y(i,l)|
 \le4^{-k-3}+\frac{2^{-k-4}+2^{-l-4}}6
 \le\frac7{192}2^{-k}<2^{-k-2}.
\]
Thus $Y$ is a regular PR name with every sample in $[6,8]$.
To verify the original equation, the internal next coordinate of
$x^{i+k+3,k+4}$ differs from $Y(i+1,k)$ by at most
$4^{-k-2}+2^{-k-4}/3$.  Adding its contribution divided by eight
to the local residual gives
\[
 \left|Y(i,k)-a_i-Y(i,k)^2/8-Y(i+1,k)/8\right|
 \le\frac{2^{-k-4}}{16}
    +\frac{4^{-k-2}+2^{-k-4}/3}{8}
 \le\frac{11}{768}2^{-k}<2^{-k-6}.
\]
This constructs a separate bounded upper solution with explicit finite
contracts.  No completed real function or unspecified square-root routine
was used to define it.  Both the bisection invariant and the finite backward
error invariant are rational inequalities suitable for quantifier-free
induction.

\subsection{The exact first-order endpoint}

Let \(\mathsf{lat}(i,k)\) be the registered PR rational code defined
by (L8). Then the main iPRA targets are the universal closures of:
\[
 \begin{gathered}
 \mathsf{inUnit}(\mathsf{lat}(i,k))=1,\\
 \mathsf{ratLe}\bigl(
   |\mathsf{lat}(i,k)-\mathsf{lat}(i,l)|,
   2^{-\min(k,l)-2}\bigr)=1,\\
 \mathsf{ratLe}\bigl(
   |\mathsf{lat}(i,k)-a_i-\mathsf{lat}(i,k)^2/8
                            -\mathsf{lat}(i+1,k)/8|,
   2^{-k-4}\bigr)=1.
 \end{gathered}
\]
Here vertical bars and fractions abbreviate the explicit rational
code operations from \cref{sec:arithmetic}. Each displayed matrix
is quantifier-free after this notation is expanded. The universal
closures are logical introductions after the finite invariant
proofs, not additional induction axioms.
